% !PRINTED X
\weberpage{X}
\inhaltpar{81}{Die Tripelsysteme der Resolventen}{310}
\inhaltpar{82}{Anwendung auf Gleichungen 8\textsuperscript{ten} Grades}{313}
\inhaltpar{83}{Metacyklische Gleichungen 8\textsuperscript{ten} Grades}{314}
\inhaltpar{84}{Biquadratische Gleichungen}{319}
\inhaltabschnitt{Elfter Abschnitt.}{Die Wendepunkte einer Curve dritter Ordnung.}
\inhaltpar{85}{Ternäre Formen und algebraische Curven}{322}
\inhaltpar{86}{Singuläre Punkte. Wendepunkte. Doppeltangenten}{324}
\inhaltpar{87}{Fundamentale Covarianten einer ternären Form}{327}
\inhaltpar{88}{Die \gesperrt{Hesse}'sche Curve}{329}
\inhaltpar{89}{Inflexionspunkte einer Curve dritter Ordnung}{331}
\inhaltpar{90}{Transformation der cubischen Form auf die canonische Form}{333}
\inhaltpar{91}{Die Invarianten der Curve dritter Ordnung und die biquadratische Gleichung}{336}
\inhaltpar{92}{Tripelgleichungen}{342}
\inhaltpar{93}{Die Gruppe der Tripelgleichungen}{344}
\inhaltpar{94}{Realitätsverhältnisse der Tripelgleichungen}{348}
\inhaltabschnitt{Zwölfter Abschnitt.}{Die Doppeltangenten einer Curve vierter Ordnung.}
\inhaltpar{95}{Anzahl der Doppeltangenten einer Curve vierter Ordnung}{351}
\inhaltpar{96}{Die \gesperrt{Steiner}'schen Complexe}{357}
\inhaltpar{97}{Complexpaare und Complextripel}{363}
\inhaltpar{98}{Die \gesperrt{Aronhold}'schen Siebener-Systeme}{367}
\inhaltpar{99}{Der \gesperrt{Hesse-Cayley}'sche Algorithmus zur Bezeichnung der Doppeltangenten}{369}
\inhaltpar{100}{Rationale Bestimmung der Curve aus einem vollständigen Siebener-Systeme}{374}
\inhaltpar{101}{Die \gesperrt{Galois}'sche Gruppe des Doppeltangentenproblems}{380}
\inhaltpar{102}{Darstellung der Gruppe}{384}
\inhaltpar{103}{Einfachheit der Gruppe des Doppeltangentenproblems}{387}
\inhaltpar{104}{Realität der Doppeltangenten}{391}
\inhaltpar{105}{Beweis der Existenz der vier Fälle}{398}
\inhaltabschnitt{Dreizehnter Abschnitt.}{Allgemeine Theorie der Gleichung fünften Grades.}
\inhaltpar{106}{Fragestellung}{403}
\inhaltpar{107}{Satz von \gesperrt{Lüroth}}{404}
\inhaltpar{108}{Resolventen mit einem Parameter}{407}
\inhaltpar{109}{Gruppe der Resolventen mit einem Parameter}{410}
\inhaltpar{110}{Die Ikosaëdergleichung}{415}
\inhaltpar{111}{Die Resolventen der Ikosaëdergleichung}{419}
\inhaltpar{112}{Die Hauptresolvente fünften Grades}{422}
